\section{Frontier shifts and separated families}\label{sec:primal-additions}

\subsection{A selected frontier shift}\label{clip:frontier}

Cumulative clipping changes the increments only where the original
test is supported. The next construction instead places a compensating
value at each relevant gate. Its two possible perturbations recover
the original tail through their difference, while the selected-endpoint
argument keeps only half of the resulting maximum excess.

\begin{proposition}[Frontier shift and selected endpoint]
\label{clip:frontier-prop}
For finite $f=P_Hf$ and $v=Q_Hv$ with $\rho(f),\rho(v)\le1$,
and every $s>0$, there is a finitely supported gate vector $a$ such that
\begin{equation}\label{clip:frontier-tests}
 \rho(a)\le1,\qquad \rho(a\pm v)\le2,\qquad
 \rho\bigl(f+s(a\pm v)\bigr)^2\le1+12s^2.
\end{equation}
For either $Z=B_\infty$ or the finite-height forest space, this
selected-endpoint method gives
\begin{equation}\label{clip:frontier-modulus}
 \widehat\delta_Z(t)\ge
 \frac{\sqrt{1+t^2/48}-1}{4}\qquad(0<t<1).
\end{equation}
\end{proposition}
\begin{proof}
Put $q_g=\rho(v\mathbf1_{U_g})$, so $\sum_gq_g^2\le1$.
At each attached gate with $q_g>0$, choose the shortest head suffix
of magnitude at least $2sq_g$, if one exists, and let $\tau_g$
be its sign. Otherwise use $\tau_g=1$. Set
$a(g)=-\tau_gq_g$ at the gates and zero elsewhere.
The gates form an antichain, so $\rho(a)^2=\sum_gq_g^2\le1$.
For $w=a\pm v$, every root prefix at $g$ has sign
$-\tau_g$ or is zero, and magnitude at most $2q_g$.
Also $\rho(w)\le2$.

In a disjoint testing family write the head and tail sums as $A_i,B_i$.
The pure tail contribution $s^2\sum_iB_i^2$ is at most $4s^2$.
A positive cross term with $|A_i|<2sq_g$ costs at most
$8s^2q_g^2$. Their total is at most $8s^2$.
For a larger positive term the selected suffix has sum $C_i$
opposite to $A_i$ and magnitude at least $2sq_g$.
Splitting it off gains
\[
 (A_i-C_i)^2+C_i^2-A_i^2
 \ge4s|A_i|q_g\ge2sA_iB_i.
\]
The simultaneous cuts are disjoint, so all these larger terms are
paid by the head norm. This proves \eqref{clip:frontier-tests}.
Unattached roots contribute only to the pure tail budget.

For a finite unit center choose its head normer $f$ and a half-norming
exterior test $v$, so $|y(v)|\ge1/2$ for $\|y\|\ge1$.
One of $a+v,a-v$ has pairing of absolute value at least $|y(v)|$,
because their difference is $2v$. Selecting that test and then the
endpoint sign gives
\[
 \max_\pm\|x\pm ty\|\ge
 \frac{1+st/2}{\sqrt{1+12s^2}}.
\]
Take $s=t/24$. The ratio equals $\sqrt{1+t^2/48}$; testing on
$f$ makes the other endpoint at least one. Finally approximate an
arbitrary unit center within one quarter of this ratio's excess.
The average retains precisely the lower bound in
\eqref{clip:frontier-modulus}.
\end{proof}

\subsection{An independent suffix threshold for root shifts}
\label{dir:primal:add:root-shifts}
A root shift can be chosen with a suffix threshold independent of the
size by which the shifted test is later multiplied. This
retains a rational estimate and also permits prefix-maximum root weights.

\begin{lemma}\label{dir:primal:add:independent-threshold}
On either the finite-height forest or the infinite tree, let
\(f=P_Hf\), \(v=Q_Hv\) be finite tests with
\(\rho(f),\rho(v)\le1\). Let finitely many \(a_g\ge0\) satisfy
\(\sum_ga_g^2\le1\) and \(|v([g,u])|\le a_g\) on every root prefix.
Fix \(\alpha>0\). For \(a_g>0\), choose the shortest attached head
suffix of magnitude at least \(\alpha a_g\), if one exists, and let
\(\tau_g\) be its sign; otherwise put \(\tau_g=1\). Strict thresholds
may be used instead. Set
\[
 a=-\sum_g\tau_ga_ge_g,\qquad h_+=a+v,\qquad h_-=a-v.
\]
For \(h\in\{a,h_+,h_-\}\) and \(0<c\le\alpha/2\),
\begin{equation}\label{dir:primal:add:independent-cost}
 \rho(h)\le2,\qquad
 \rho(f+ch)^2\le1+4c\alpha+4c^2.
\end{equation}
Both the full cone norms \(q_g\) and the prefix maxima \(m_g\) may be
used for \(a_g\).
\end{lemma}
\begin{proof}
Since the gates form an antichain,
\(\rho(a)^2=\sum_ga_g^2\le1\). Each root prefix of \(h\) has sign
\(-\tau_g\) or is zero and magnitude at most \(2a_g\), and
\(\rho(h)\le2\). Zero weights give zero cross terms; unattached roots have no head crossing.

For a disjoint test family write \(A_i,B_i\) for head and tail sums and
\(L=\sum_iA_i^2\le1\). Small positive products, with
\(|A_i|\le\alpha a_g\), have total at most \(2\alpha\).
For every other positive product, the simultaneous suffix cuts in
\eqref{clip:split-gain} release at least \(2|A_i|\alpha a_g\).
Their total release is at most \(1-L\), whereas \(|B_i|\le2a_g\).
Thus the large positive products total at most \((1-L)/\alpha\).
Since \(\sum_iB_i^2\le4\),
\[
 \sum_i(A_i+cB_i)^2
 \le L+2c\bigl(2\alpha+(1-L)/\alpha\bigr)+4c^2
 \le1+4c\alpha+4c^2.
\]
This uses precisely \(2c\le\alpha\). Strict thresholds use the same
small class, including equality. Formula~\eqref{clip:budgets} verifies
both proposed choices of weights.
\end{proof}

The next two endpoint estimates concern the finite-height forest dual $X$.

\paragraph{The rational gap.}\label{dir:primal:add:rational}
At a finite unit center take its head normer \(f\), and a finite tail
normer \(v\) with \(\rho(v)=1\) and \(y(v)>1/2\), for unit
\(y\in\ker P_H\). Use \(a_g=q_g\), so \(\sum_gq_g^2=1\).
The difference \(h_+-h_-=2v\) gives one \(|y(h)|>1/2\).
The mixed head--tail pairings vanish, and
\[
 \max_{\pm}\|x\pm ty\|
 \ge\frac{1+ct/2}{\sqrt{1+4c\alpha+4c^2}}
 \ge\frac{1+ct/2}{1+2c\alpha+2c^2}.
\]
Choose \(\alpha=t/16\), \(c=t/32\), for any \(t>0\). The ratio equals
\(1+d_0(t)\), where
\[
 d_0(t)=\frac{5t^2}{512+3t^2}.
\]
Approximation of any unit center by a finite unit center within
\(d_0(t)/2\) retains gain \(d_0(t)/2\), with the same finite-coordinate
annihilator.

\paragraph{Prefix weights at the actual center.}
\label{dir:primal:add:root-actual}
Fix \(t>0\), \(0<c<t/6\), and put
\(L_c=\sqrt{1+12c^2}\), \(b=1+ct-L_c>0\).
Choose
\[
 0<\eta<\min\{1,b/[2(1+ct+2c)]\},\qquad
 \gamma=b/(2L_c).
\]
For any unit \(x\), choose a finite ancestral \(H\) with
\(\|Q_Hx\|<\eta\) and a head test \(f\) with
\(\rho(f)\le1\), \(x(f)\ge1-\eta\).
For a unit \(y\in\ker P_H\), choose a finite tail test \(v\) with
\(\rho(v)\le1\), \(y(v)>1-\eta\). In
Lemma~\ref{dir:primal:add:independent-threshold}, use prefix weights
\(a_g=m_g\), strict threshold \(\alpha=2c\), and the deepest qualifying
starting vertex, which is the shortest qualifying suffix.
One \(h\in\{h_+,h_-\}\) has \(|y(h)|>1-\eta\), while
\(|x(h)|<2\eta\). Consequently
\begin{equation}\label{dir:primal:add:root-actual-gain}
 \max_{\pm}\|x\pm ty\|
 \ge\frac{1-\eta+c((1-\eta)t-2\eta)}{\sqrt{1+12c^2}}
 >1+\gamma.
\end{equation}
Indeed \(L_c\le1+6c^2<1+ct\), and the numerator is
\(L_c+b-\eta(1+ct+2c)>L_c+b/2\).

\subsection{Two tests for a bounded separated family}
\label{dir:primal:add:fixed-bound}
The final variant uses a tail witness selected from an entire separated
family. Its member need not annihilate the head. Testing the two endpoints
with different primal vectors cancels that unknown head evaluation.

\begin{proposition}\label{dir:primal:add:separated}
For each \(D\ge1\), there exists \(c_D>0\) such that, if
\(x,z_i\in X\), \(1\le\|x\|\le D\), \(\|z_i\|\le2D\), and
\(\|z_i-z_j\|\ge1\) for distinct indices in an infinite sequence, then
some \(i\) satisfies \(\max_{\pm}\|x\pm z_i\|>\|x\|+c_D\).
\end{proposition}
\begin{proof}
First choose \(c>0\), depending only on \(D\), with
\(D(\sqrt{1+16c^2}-1)<c/16\), and then choose
\(0<\eta<c/[64(1+c)]\).
Choose a head \(H\), a head test \(f\), and a center tail error so that
\(\rho(f)\le1\), \(x(f)>\|x\|-\eta\), and \(\|Q_Hx\|<\eta\).
Boundedness in the finite-dimensional head gives distinct indices with
\(\|P_H(z_i-z_j)\|<1/4\), hence their tail difference has norm greater
than \(3/4\). A finite tail test \(v\), \(\rho(v)\le1\), has
\(|(z_i-z_j)(v)|>1/2\). Rename one index so \(|z_i(v)|>1/4\).

Use prefix weights \(a_g=m_g\) and strict threshold \(2c m_g\) in
Lemma~\ref{dir:primal:add:independent-threshold}. The two perturbations
\(h_0=a\), \(h_1=a+v\) differ by \(v\), have norms at most one and two,
and have the same inward prefix signs with magnitudes at most \(2m_g\).
One choice \(h\) has \(|z_i(h)|>1/8\).
We retain the following direct whole-crossing estimate:
\begin{equation}\label{dir:primal:add:sixteen}
 \rho(f+ch)^2\le1+16c^2.
\end{equation}
For a crossing write its entire sum as \(A+B\), where \(B\) already
includes the factor \(c\), so \(|B|\le2cm_g\).
If \(|A|\le2cm_g\), its square is at most \(A^2+12c^2m_g^2\).
If \(|A|>2cm_g\) and \(AB\le0\), it is at most \(A^2\).
Otherwise the chosen suffix has sum \(C\) opposite to \(A\), with
\(|C|>2cm_g\ge|B|\), and
\[
 (A-C)^2+C^2
 =A^2+2|A||C|+2C^2\ge(A+B)^2.
\]
The simultaneous head cuts remain disjoint. Crossings and wholly head
segments therefore cost at most \(1+12c^2\); wholly tail segments add
at most \(c^2\rho(h)^2\le4c^2\). This proves
\eqref{dir:primal:add:sixteen}.

Put \(M=\max_{\pm}\|x\pm z_i\|\), and choose \(\epsilon\in\{-1,1\}\)
so that \(\epsilon z_i(h)=|z_i(h)|\). Test \(x+\epsilon z_i\) on
\(f+ch\), and \(x-\epsilon z_i\) on \(f\), and add. The head evaluations
of \(z_i\) cancel. With \(L=\sqrt{1+16c^2}\), this gives
\begin{equation}\label{dir:primal:add:two-test}
 M(1+L)\ge2(\|x\|-\eta)+c(1/8-2\eta).
\end{equation}
Here \(x(h)>-2\eta\). It follows that
\[
 (1+L)(M-\|x\|)
 >c/16-2\eta(1+c)>c/32.
\]
Thus \(c_D=c/[32(1+L)]\) works. The parameter choice depends on the fixed
bound \(D\); no unit-tail hypothesis on the selected \(z_i\) was used.
\end{proof}
